\subsection{Variational energies and feasible spaces}
\label{sec:models}

The three \(p\)-energies are defined in \cref{eq:pure-energy,eq:reg-energy,eq:reaction-energy}. The
remaining models have the following energies, written for \(d\leq2\):
\begin{align}
 \E_{\rm lin}(u)&=\int_\Omega\left(\frac12\abs{\nabla u}^2+
 \frac12u^2\right)\dd x-\ip fu,
 \label{eq:linear-energy}\\
 \E_{\rm cub}(u)&=\int_\Omega\left(\frac12\abs{\nabla u}^2+
 \frac14u^4\right)\dd x-\ip fu,
 \label{eq:cubic-energy}\\
 \E_{\sinh}(u)&=\int_\Omega\left(\frac12\abs{\nabla u}^2+
 \cosh u-1\right)\dd x-\ip fu.
 \label{eq:sinh-energy}
\end{align}
Here \(\ip fu\) denotes the duality pairing between the relevant
energy space and its dual. When the load admits an integral
representation, \(\ip fu=\int_\Omega fu\,\mathrm dx\).
Subtracting \(1\) from \(\cosh u\) normalizes the potential to zero
at \(u=0\) and leaves the first variation, hence the
Euler--Lagrange equation, unchanged.

\begin{assumption}[Data and feasible spaces for the six models]
\label[assumption]{ass:model-data}
The following model-specific conditions are imposed.
\begin{enumerate}[label=\textup{(\roman*)}]
 \item For the linear and cubic models, \(f\in (H^1(\Omega))^*\). The sinh energy has the same load space and is treated as an extended-real convex integral functional with nonempty effective domain. Its local derivative and Hessian estimates are formulated on an \(L^\infty\)-bounded neighborhood.
 \item For the pure and regularized \(p\)-energies, \(f\in X_{p,\circ}^*\), or equivalently the Neumann
 load annihilates constants before restriction to \(X_{p,\circ}\).
 \item For the reaction \(p\)-energy, \(f\in(W^{1,p})^*\). In the regularized model, the parameter
 \(\eps\) in \cref{eq:reg-energy} is fixed and positive.
\end{enumerate}
\end{assumption}

The local Hessian structure of each model is needed for the transfer estimates of
\Cref{sec:c5-hessian}. The bilinear forms at the exact solution are
\begin{align*}
 a_{\rm lin}(v,w)&=\int_\Omega(\nabla v\cdot\nabla w+vw)\dd x,\\
 a_{\rm cub,*}(v,w)&=\int_\Omega(\nabla v\cdot\nabla w+3(u^*)^2vw)\dd x,\\
 a_{\sinh,*}(v,w)&=\int_\Omega(\nabla v\cdot\nabla w+\cosh(u^*)vw)\dd x.
\end{align*}

\begin{proposition}[Well-posedness]
\label{prop:wellposedness}
Under \Cref{ass:model-data}, \cref{eq:linear-energy,eq:cubic-energy,eq:sinh-energy} have unique minimizers
in \(H^1(\Omega)\); \cref{eq:pure-energy,eq:reg-energy} have unique minimizers in \(X_{p,\circ}\); and
\cref{eq:reaction-energy} has a unique minimizer in \(W^{1,p}(\Omega)\).
\end{proposition}

\begin{proof}[Proof sketch]
The quadratic, strictly convex, and $p$-growth terms give coercivity on the stated feasible spaces; strict convexity gives uniqueness.  The pure Neumann nullspace is removed by the zero-mean restriction.
\end{proof}

For a \(p\)-diffusion density, the Hessian matrix in gradient variables is
\[
 \abs\xi^{p-2}I+(p-2)\abs\xi^{p-4}\xi\otimes\xi
\]
at nonzero \(\xi\) in the pure case, with \(\abs\xi^2\) replaced by
\(\eps^2+\abs\xi^2\) in the regularized case. The local Hessian statements that follow concern
\(p\geq2\).

\begin{proposition}[Sufficient continuous local Hessian conditions]
\label{prop:model-hessian}
In dimensions \(d\leq2\), the cubic Hessian is locally Lipschitz as a bilinear form on \(H^1\) and is coercive at any solution \(u^*\not\equiv0\). The sinh Hessian is coercive and locally Lipschitz on an \(L^\infty\)-bounded neighborhood, with state perturbations measured in \(H^1\). For fixed \(\eps>0\), the regularized \(p\)-Hessian is uniformly elliptic with lower bound proportional to \(\eps^{p-2}\), and locally Lipschitz on bounded subsets of \(W^{1,\infty}\).
\end{proposition}

\begin{proof}[Proof sketch]
In dimensions $d\leq2$, Sobolev embedding controls the cubic coefficient on bounded $H^1$ neighborhoods, while an $L^\infty$ bound controls the sinh coefficient.  Differentiating the regularized radial flux gives uniform ellipticity and bounded derivative on bounded-gradient sets.
\end{proof}

\begin{proposition}[Restricted stability for the regularized \(p\)-energy]
\label{prop:regularized-restricted}
Let \(Y_M\subset X_{p,\circ}\cap W^{1,\infty}(\Omega)\) be finite dimensional, equipped with a norm \(\norm\cdot_{0,M}\), and let \(U_M\subset Y_M\) be a bounded convex set. For fixed \(p\ge2\) and \(\eps>0\), there are constants \(\mu_M>0\) and \(L_M<\infty\) such that, for \(u,z\in U_M\) and \(v,w\in Y_M\),
\[
 \begin{aligned}
 \E''(u)[v,v]&\ge\mu_M\norm v_{0,M}^2,\\
 |(\E''(u)-\E''(z))[v,w]|
 &\le L_M\norm{u-z}_{0,M}\norm v_{0,M}\norm w_{0,M}.
 \end{aligned}
\]
The constants depend on \(U_M\), \(Y_M\), the reference norm, \(p\), and \(\eps\).
\end{proposition}

\begin{proof}[Proof sketch]
Finite-dimensional norm equivalence gives a Hessian-Lipschitz constant on the bounded set. Uniform ellipticity and the zero-mean condition give coercivity in the reference norm.
\end{proof}

The linear, cubic, and sinh models are treated through their quadratic or local Hessian structure. For the regularized \(p\)-energy, uniform ellipticity gives the stated finite-dimensional estimates. The pure and reaction \(p\)-energies are described by the global energy--distance relations and the finite-trajectory projection analysis.
